\scanpage{98}
On aura alors
\[
\left\{\begin{aligned}
&\deleted{\unsure{\kappa_{i'}=|\omega_{i'}=}}\quad\omega_{i'}=-\omega_i\\
&\text{i.e. }\ i\longmapsto\omega_i\quad\text{est une bijection }\deleted{\unsure{\text{des}|\text{de}}}\\
&\hspace{4em}\alpha=\alpha_\kappa:\mu_4^*\xrightarrow{\sim}\underline\omega
\end{aligned}\right. \tag*{(68)}
\]
% slanted
\margin{\uline{NB} Quand on se donne d'avance l'iso. $\alpha:\mu_4^\times\simeq\underline\omega$, alors les deux relations (67) sont équivalentes \ldots}
On va expliciter la donnée d'un $\kappa$, \uline{associé} à un isom. \deleted{$\alpha:\mu_4^\times\simeq\underline\omega$,} \inserted{donné au préalable} \deleted{\unsure{$\alpha$|à}} $\alpha=(i\longrightarrow\omega_i):\mu_4^\times\simeq\underline\omega$, via les $\check\kappa_i$. On aura, pour $\delta\in\underline t^*$
\[
\check\kappa_i(\check\delta)=\underset{\in F_t}{\underline{\psi_i(\delta)}}\mathbin{\wedge}\omega_i
=\underline{\psi^{\omega_i}(\delta)}\mathbin{\wedge}\omega_i\,, \tag*{-(69)}
\]
où
\[
\bigl(\psi^\omega(\delta)\bigr)_{\omega\in\underline\omega}\ \in (F_t)^{\underline\omega} \tag*{(70)}
\]
est une famille de \deleted{\unsure{éléments|élément}} deux éléments de $F_t$, indexée par $\underline\omega$. La condition\deleted{s} \inserted{(65) ou} (66) s'explicite\deleted{\unsure{nt}} comme
\[
\underbrace{\psi^{\omega_{i'}}(\delta)\mathbin{\wedge}\omega_{i'}}_{
\substack{\psi^{\omega'}(\delta)\mathbin{\wedge}\omega'\\
=-\psi^{\omega'}(\delta)\mathbin{\wedge}\omega\\
\text{si }\deleted{\unsure{\omega'\text{ pour}}}\ \text{\unsure{pos|pose|on pose}}\\
\omega\overset{\mathrm{def}}{=}\omega_i\\
\omega'=\text{autre él. de }\underline\omega}}
=\frac1i\underbrace{\psi^{\omega_i}(\delta)\mathbin{\wedge}\omega_i}_{\psi^\omega(\delta)\mathbin{\wedge}\omega}
\]
i.e. $\psi^{\omega'}(\delta)=\underbrace{-\dfrac1i}_{i}\psi^\omega(\delta)$

i.e
\[
\psi^{\omega'}(\delta)=i_\omega\psi^\omega(\delta)
\qquad\text{où }\ i_\omega\overset{\mathrm{def}}{=}\alpha^{-1}(\omega) \tag*{(71)}
\]
On doit avoir de plus
\[
\psi^\omega(\lambda\delta)=\frac1\lambda\psi^\omega(\delta)\,, \tag*{(72)}
\]
\deleted{qui ex} et enfin il faut exprimer (67),
% archive sheet number: 97
